% !TeX root = ../../tesis.tex
%%%%%%%%%%%%%%%%%%%%%%%%%%%%%%%%%%%%%%%%%%%%%%%%%%%%%%%%%%%%%%%%%%%
% §2.1.2 HERRAMIENTAS COMPUTACIONALES Y DATOS ABIERTOS PARA LA PLANEACIÓN
%        — EN BLANCO
% Origen 1: 3-Conceptos-Indicadores/3-1-Definiciones/3-1-3-Planeacion.tex
%   (118 l.; §3.1.3 «Planeación Estratégica» completo, pedido del asesor):
%   - l. 8-44  «Caracterización Dinámica de la Oferta de Transporte»
%              (incluye fig:gtfs-estructura y su cita)
%   - l. 45-72 «Toma de Decisiones Basada en Algoritmos y Comparación de Redes»
%   - l. 73-118 «Enfoque en la Sostenibilidad y el Cambio Modal»
%   Autores: fortin2016innovative, moreno2023, ugur2025, lgmsv2022,
%   hidalgo2013brt, pojani2015sustainable.
% Origen 2: 2-MarcoTeorico/2-5-SIG-Datos-Abiertos.tex (103 l.):
%   - l. 4-20  entrada de la sección
%   - l. 21-50 «VGI y cobertura de datos a escala metropolitana»
%   - l. 51-77 «Datos abiertos, comunidad y equidad de acceso»
%   Autores: haklay2010good, barrington2017world, neis2012street,
%   gurstein2011open.
%   (l. 78-103 «De los datos abiertos a la herramienta replicable» NO va
%   aquí: pasa al final de 2-1-5-Sintesis, decisión D36-6.)
% Al mover: texto continuo sin \subsubsection; organizar por tema
%   (grafos+GTFS → software propio → datos abiertos/VGI → brecha técnica).
%%%%%%%%%%%%%%%%%%%%%%%%%%%%%%%%%%%%%%%%%%%%%%%%%%%%%%%%%%%%%%%%%%%
